\subsection{Behler \& Parrinello (2007): high-dimensional NN potential for bulk silicon}
\label{sec:appendix-example-behler2007}

\paragraph{Q1 -- Bibliographic identification.}
Behler and Parrinello introduce the foundational
high-dimensional neural-network potential (HDNNP) construction:
total energy as a sum of per-atom subnet contributions over
shared-weight feed-forward NNs that take atom-centred symmetry
functions as input. The method is demonstrated on bulk silicon.
Published in \emph{Phys.\ Rev.\ Lett.} 98, 146401 (2007).
\lstinputlisting[language=turtle]{artifacts/kg/listings/behler2007/q01-bibliographic.ttl}

\paragraph{Q2 -- Material system.}
The demonstration system is bulk silicon. The training set covers
the diamond-cubic semiconducting phase, the high-pressure phases
(including $\beta$-tin), and the liquid metallic phase, so we
record \texttt{Si} as a single $\MaterialSystemC$ at the
elemental level with the phase coverage captured in
$\dlRole{materialClass}$.
\lstinputlisting[language=turtle]{artifacts/kg/listings/behler2007/q02-material.ttl}

\paragraph{Q3 -- Reference calculation method.}
Reference data are produced from DFT calculations in the local
density approximation (LDA).
\lstinputlisting[language=turtle]{artifacts/kg/listings/behler2007/q03-reference-method-type.ttl}

\paragraph{Q4 -- Reference settings.}
DFT calculations are carried out with PWscf using a 20~Ry
($\approx 272$~eV) plane-wave cutoff in combination with a
Vanderbilt ultrasoft pseudopotential, a $3 \times 3 \times 3$
k-point mesh, and Fermi smearing of 0.1~eV (to aid convergence in
the metallic phases). Frozen-core treatment beyond the
ultrasoft-pseudopotential partition is not explicitly reported.
\lstinputlisting[language=turtle]{artifacts/kg/listings/behler2007/q04-reference-settings.ttl}

\paragraph{Q5 -- Training dataset.}
Approximately 9{,}000 DFT energies are computed in total: 8{,}200
are used for training the NN and 800 form a held-out test set.
The dataset is in-house and covers energies and atomic forces.
\lstinputlisting[language=turtle]{artifacts/kg/listings/behler2007/q05-dataset.ttl}

\paragraph{Q6 -- Sampling strategies.}
Three strategies are used: enumeration of crystal structures
(including high-pressure polymorphs), DFT MD snapshots at
different pressures and temperatures, and a self-consistent
iterative refinement loop in which best fits drive MD,
hybrid Monte Carlo, and metadynamics runs whose representative
configurations are recalculated with DFT and added to the
training set when their RMSE exceeds the current fit error---a
precursor to today's active-learning workflows.
\lstinputlisting[language=turtle]{artifacts/kg/listings/behler2007/q06-sampling.ttl}

\paragraph{Q7 -- MLIP method, components, implementation.}
The method is the original Behler--Parrinello high-dimensional
NN potential: the total energy decomposes as
$E = \sum_i E_i$, with each $E_i$ produced by a shared-weight
feed-forward subnet whose inputs are atom-centred symmetry
functions (radial Gaussian sums and angular triplet sums). The
training loss is the mean squared error on the per-configuration
total DFT energy. The implementation is an in-house extension of
the Lorenz--Gro\ss{}--Scheffler NN code; the training algorithm
is not explicitly named.
\lstinputlisting[language=turtle]{artifacts/kg/listings/behler2007/q07-method.ttl}

\paragraph{Q8 -- Hyperparameter settings.}
Reported hyperparameter settings: a cutoff radius of 6~\AA{},
48 symmetry functions per atom, and subnets with 2 hidden layers
of $\approx 40$ nodes each (a few thousand fitting parameters in
total). The hyperbolic tangent is the hidden-layer activation;
the output layer uses a linear activation. Numerical Gaussian
parameters $\eta, R_s, \zeta$ are not enumerated.
\lstinputlisting[language=turtle]{artifacts/kg/listings/behler2007/q08-hyperparameter-settings.ttl}

\paragraph{Q9 -- MLIP run and trained model.}
A single training run produces one HDNNP for bulk silicon. The
self-consistent iterative refinement is recorded as a sampling
strategy in Q6 rather than as a multi-stage run.
\lstinputlisting[language=turtle]{artifacts/kg/listings/behler2007/q09-run-and-model.ttl}

\paragraph{Q10 -- Benchmark results.}
Headline accuracy figures are an energy RMSE of 4--5~meV/atom
on the optimisation set and 5--6~meV/atom on the independent test
set, plus a force-component RMSE of $\approx 0.2$~eV/\AA{} (vs.\
DFT-LDA). Downstream qualitative checks include reproducing the
DFT energy--volume curves and transition pressures across the
silicon polymorphs, and the radial distribution function of the
3000~K silicon melt.
\lstinputlisting[language=turtle]{artifacts/kg/listings/behler2007/q10-benchmark-results.ttl}

\paragraph{Q11 -- Computational resources.}
\emph{Not reported.} The paper gives no training duration, GPU
hours, training hardware, peak memory, or per-atom inference
time. They do quote a relative inference speed of $\sim$5 orders
of magnitude faster than DFT for a 64-atom cell---a comparative
figure rather than an ontology-recordable
$\dlRole{inferenceTimePerAtom}$ value.
\lstinputlisting[language=turtle]{artifacts/kg/listings/behler2007/q11-resources.ttl}

\paragraph{Q12 -- Gaps.}
Beyond Q11 (compute resources entirely absent), the following
ontology-expressible items are not reported: the training
algorithm or optimiser used to fit the NN weights;
software-implementation versioning (this is a precursor to
today's named HDNNP packages such as RuNNer or n2p2); explicit
$\dlRole{frozenCore}$ treatment beyond the ultrasoft-pseudopotential
partition; numerical values of the symmetry-function parameters
$\eta$, $R_s$, $\zeta$; and a controlled-vocabulary individual
for the self-consistent iterative refinement (encoded as an
ad-hoc \texttt{rdfs:label}/\texttt{rdfs:comment} sampling
instance, since active learning post-dates this paper).
